\documentclass[a4paper]{article}

\usepackage[spanish]{babel}
\usepackage{graphicx}
\usepackage{amsmath, amssymb, mathrsfs}
\usepackage[margin=2cm]{geometry}
\usepackage{fancyhdr}
\usepackage{enumerate}
\usepackage[shortlabels]{enumitem}
\usepackage{parskip}
\usepackage[most]{tcolorbox}
\usepackage[hidelinks]{hyperref}
\usepackage{float}
\usepackage{booktabs}



% cabecera
\pagestyle{fancy}
\fancyhead[l]{Física Matemática I}
\fancyhead[c]{Átomo de Hidrógeno}

\renewcommand{\headrulewidth}{0.2pt} % linea horizontal

\begin{document}

\begin{titlepage}
  \centering
  {\Large Universidad de Antioquia\par}
  \vspace{2cm}
  {\huge\bfseries Física Matemática I\par}
  \vspace{0.4cm}
  {\Large Taller del átomo de hidrógeno \par}
  \vspace{2.5cm}
  {\large Autor:\par}
  \vspace{0.2cm}
  {\large Daniel Soto Villada\par}
  {\large Estudiante de Astronomía\par}
  \vfill
  {\large \today\par}
\end{titlepage}

\newpage
\tableofcontents
\newpage


\section{Aplicación: Átomo de Hidrógeno}

En mecánica cuántica el comportamiento del electrón en un átomo de Hidrógeno está descrito por la función de onda $\psi$ que se obtiene al solucionar la ecuación de Schrodinger

$$
-\frac{\hbar^2}{2m}\nabla^2\psi -\frac{K}{r}\psi = E\psi
$$

donde $K = \frac{e^2}{4\pi\epsilon_0}$

\subsection{a) Construcción de la ecuación radial}
Usando la separación de variables $\psi(r,\theta,\varphi) = \frac{R(r)}{r}Y(\theta,\varphi)$, muestre que la ecuación diferencial radial está dada por:

$$
\ddot{R}(\alpha x) + \left[ -\frac{l(l+1)}{x^2} + \frac{\alpha K 2m}{x\hbar^2} + \frac{2mE}{\hbar^2} \right] R(\alpha x) = 0 \,, \eqno{(44)}
$$

donde $r = \alpha x$.

\subsubsection{Solución}
Llevamos la ecuación a la forma $(\nabla^2 + \text{factor})\psi = 0$ y nos queda

$$\left[ \nabla^2 + \frac{2m}{\hbar^2}\left(E + \frac{K}{r}\right) \right]\psi = 0$$

Veamos la expresión del laplaciano en coordenadas esféricas.

$$\nabla^2 = \frac{1}{r^2} \partial_r (r^2 \partial_r) + \frac{1}{r^2 \sin\theta} \partial_\theta (\sin\theta \partial_\theta) + \frac{1}{r^2 \sin^2\theta} \partial_\varphi^2$$

La aplicamos a $\psi$

$$\nabla^2 \left[ \frac{R}{r}Y \right] = Y \cdot \frac{1}{r^2} \partial_r \left( r^2 \partial_r \left(\frac{R}{r}\right)\right) + \frac{R}{r} \cdot \left[ \frac{1}{r^2 \sin\theta} \partial_\theta (\sin\theta \partial_\theta Y) + \frac{1}{r^2 \sin^2\theta} \partial_\varphi^2 Y \right]$$

$$\partial_r \left(\frac{R}{r}\right) = \frac{\mathrm{d}}{\mathrm{d}r}\left(\frac{R}{r}\right) = \frac{R'}{r} - \frac{R}{r^2}$$

$$r^2 \partial_r \left(\frac{R}{r}\right) = r^2 \left( \frac{R'}{r} - \frac{R}{r^2} \right) = r R' - R$$

$$\partial_r \left( r R' - R \right) = \frac{\mathrm{d}}{\mathrm{d}r} (r R' - R) = R' + r R'' - R' = r R''$$

$$\frac{1}{r^2} \partial_r \left( r^2 \partial_r \left(\frac{R}{r}\right)\right) = \frac{1}{r^2} (r R'') = \frac{R''}{r}$$


Introducimos este término radial simplificado en la ecuación de Schrödinger original:

$$\left[ \nabla^2 + \frac{2m}{\hbar^2}\left(E + \frac{K}{r}\right) \right]\psi = 0$$

$$Y \left( \frac{R''}{r} \right) + \frac{R}{r} \left[ \frac{1}{r^2 \sin\theta} \partial_\theta (\sin\theta \partial_\theta Y) + \frac{1}{r^2 \sin^2\theta} \partial_\varphi^2 Y \right] + \frac{2m}{\hbar^2}\left(E + \frac{K}{r}\right) \frac{R}{r} Y = 0$$

Pasamos la parte radial a la izquierda y la angular a la derecha

$$Y \left( \frac{R''}{r} \right) + \frac{2m}{\hbar^2}\left(E + \frac{K}{r}\right) \frac{R}{r} Y = -\frac{R}{r} \left[ \frac{1}{r^2 \sin\theta} \partial_\theta (\sin\theta \partial_\theta Y) + \frac{1}{r^2 \sin^2\theta} \partial_\varphi^2 Y \right]$$

Para aislar por completo las variables de cada lado, multiplicamos toda la ecuación por el factor $\frac{r^3}{R Y}$:

$$\frac{r^3}{R Y} \cdot Y \left( \frac{R''}{r} \right) + \frac{r^3}{R Y} \cdot \frac{2m}{\hbar^2}\left(E + \frac{K}{r}\right) \frac{R}{r} Y = \frac{r^2 R''}{R} + \frac{2m r^2}{\hbar^2}\left(E + \frac{K}{r}\right)$$

$$\frac{r^3}{R Y} \cdot \left( -\frac{R}{r} \right) \left[ \frac{1}{r^2 \sin\theta} \partial_\theta (\sin\theta \partial_\theta Y) + \frac{1}{r^2 \sin^2\theta} \partial_\varphi^2 Y \right] = -\frac{r^2}{Y} \left[ \frac{1}{r^2 \sin\theta} \partial_\theta (\sin\theta \partial_\theta Y) + \frac{1}{r^2 \sin^2\theta} \partial_\varphi^2 Y \right]$$

$$\frac{r^2 R''}{R} + \frac{2m r^2}{\hbar^2}\left(E + \frac{K}{r}\right) = -\frac{1}{Y} \left[ \frac{1}{\sin\theta} \partial_\theta (\sin\theta \partial_\theta Y) + \frac{1}{\sin^2\theta} \partial_\varphi^2 Y \right]$$


Vamos a construir la ecuación radial

Igualamos la parte radial a la constante de separación $\ell(\ell+1)$:

$$\frac{r^2 R''}{R} + \frac{2m r^2}{\hbar^2}\left(E + \frac{K}{r}\right) = \ell(\ell+1)$$

$$R'' + \frac{2m}{\hbar^2}\left(E + \frac{K}{r}\right)R = \frac{\ell(\ell+1)}{r^2}R$$

$$R'' + \frac{2m}{\hbar^2}\left(E + \frac{K}{r}\right)R - \frac{\ell(\ell+1)}{r^2}R = 0$$

$$R'' + \left[ \frac{2mE}{\hbar^2} + \frac{2mK}{\hbar^2 r} - \frac{\ell(\ell+1)}{r^2} \right] R = 0$$

Vamos a aplicar el cambio de variable.

$$r = \alpha x \implies x = \frac{r}{\alpha}$$

Calculamos la derivada de $x$ respecto a $r$:

$$\frac{dx}{dr} = \frac{1}{\alpha}$$

Aplicamos la regla de la cadena

$$R' = \frac{dR}{dr} = \frac{dR}{dx}\frac{dx}{dr}$$

Sustituyendo el valor de $\frac{dx}{dr}$ y utilizando la notación donde la derivada respecto a $x$ se denota con un punto, es decir, $\frac{dR}{dx} = \dot{R}(\alpha x)$, obtenemos:

$$R' = \frac{1}{\alpha} \dot{R}(\alpha x)$$

$$R'' = \frac{d}{dr}\left(\frac{dR}{dr}\right) = \frac{d}{dx}\left(\frac{dR}{dr}\right) \frac{dx}{dr}$$

$$R'' = \frac{d}{dx}\left(\frac{1}{\alpha} \dot{R}(\alpha x)\right) \cdot \left(\frac{1}{\alpha}\right)$$

$$R'' = \frac{1}{\alpha^2} \ddot{R}(\alpha x)$$

Tengamos a la mano la ecuación diferencial donde $R=R(r)$

$$R'' + \left[ \frac{2mE}{\hbar^2} + \frac{2mK}{\hbar^2 r} - \frac{\ell(\ell+1)}{r^2} \right] R = 0$$

Sustituimos $R''$, $r = \alpha x$, $r^2 = \alpha^2 x^2$ y escribimos $R$ como función explícita de $\alpha x$:

$$\frac{1}{\alpha^2} \ddot{R}(\alpha x) + \left[ \frac{2mE}{\hbar^2} + \frac{2mK}{\hbar^2 \alpha x} - \frac{\ell(\ell+1)}{\alpha^2 x^2} \right] R(\alpha x) = 0$$

$$\ddot{R}(\alpha x) + \left[ \frac{2mE\alpha^2}{\hbar^2} + \frac{2mK\alpha}{\hbar^2 x} - \frac{\ell(\ell+1)}{x^2} \right] R(\alpha x) = 0$$

$$\ddot{R}(\alpha x) + \left[ -\frac{\ell(\ell+1)}{x^2} + \frac{2mK\alpha}{\hbar^2 x} + \frac{2mE\alpha^2}{\hbar^2} \right] R(\alpha x) = 0$$



\subsection{b)}
Comprando la familia de ecuaciones de Laguerre, muestre que:

$$
E_n = -\frac{m}{2\hbar^2}\frac{\tilde{k}^2}{n^2} \hspace{1.5cm} n = p + l + 1 = \{1, 2, 3, \cdots \}
$$

La familia de ecuaciones de Laguerre se escriben así:

$$\ddot{\psi}_p^k(x) + \psi_p^k(x) \left[ -\frac{(k^2 - 1)}{4x^2} + \frac{2p + k + 1}{2x} - \frac{1}{4} \right] = 0$$

La familia de ecuaciones de Laguerre nos ayuda a obtener la solución para la función de onda (debidamente normalizada).

Aquí tienes la transcripción aplicando la notación solicitada:

$$\psi_{n\ell m}(r, \theta, \varphi) = \sqrt{ \left( \frac{2}{na_0} \right)^3 \frac{(n-\ell-1)!}{2n(n+\ell)!} } e^{-\rho/2} \rho^\ell L_{n-\ell-1}^{2\ell+1}(\rho) Y_{\ell m}(\theta, \varphi) ,$$

donde $\rho = 2r/na_0$, $a_0 = \frac{4\pi\epsilon_0\hbar^2}{me^2}$ es el radio de Bohr, $n \geq 1$, $0 \leq \ell < n$, $|m| \leq \ell$.

\subsubsection{Solución}

Veamos las ecuaciones que tenemos que igualar:

$$\ddot{R}(\alpha x) + \left[ -\frac{\ell(\ell+1)}{x^2} + \frac{2mK\alpha}{\hbar^2 x} + \frac{2mE\alpha^2}{\hbar^2} \right] R(\alpha x) = 0$$

$$\ddot{\psi}_p^k(x) + \psi_p^k(x) \left[ -\frac{(k^2 - 1)}{4x^2} + \frac{2p + k + 1}{2x} - \frac{1}{4} \right] = 0$$

Igualamos los coeficientes que dividen por $x^2$:

$$-\ell(\ell+1) = -\frac{k^2 - 1}{4}$$

$$4\ell^2 + 4\ell = k^2 - 1$$

$$4\ell^2 + 4\ell + 1 = k^2$$

El lado izquierdo es un trinomio cuadrado perfecto, lo que nos permite simplificar:

$$(2\ell + 1)^2 = k^2 \implies k = 2\ell + 1$$

Ahora, igualamos los coeficientes que dividen por $x$:

$$\frac{2mK\alpha}{\hbar^2} = \frac{2p + k + 1}{2}$$

Sustituimos el valor de $k = 2\ell + 1$ que acabamos de encontrar:

$$\frac{2mK\alpha}{\hbar^2} = \frac{2p + (2\ell + 1) + 1}{2}$$

$$\frac{2mK\alpha}{\hbar^2} = \frac{2p + 2\ell + 2}{2} = p + \ell + 1$$

Definimos el número cuántico principal como $n = p + \ell + 1$.

$$\frac{2mK\alpha}{\hbar^2} = n$$

De aquí, despejamos el factor de escala $\alpha$:

$$\alpha = \frac{n\hbar^2}{2mK}$$

Finalmente, igualamos los términos que no dependen de $x$:

$$\frac{2mE\alpha^2}{\hbar^2} = -\frac{1}{4}$$

Sustituimos la expresión de $\alpha$ que obtuvimos en el paso anterior:

$$\frac{2mE}{\hbar^2} \left( \frac{n\hbar^2}{2mK} \right)^2 = -\frac{1}{4}$$

Expandimos el cuadrado y simplificamos:

$$\frac{2mE}{\hbar^2} \left( \frac{n^2\hbar^4}{4m^2K^2} \right) = -\frac{1}{4}$$

$$E \left( \frac{n^2\hbar^2}{2mK^2} \right) = -\frac{1}{4}$$

Por último, despejamos la energía $E$ (que ahora llamaremos $E_n$ por su dependencia de $n$):

$$E_n = -\frac{1}{4} \left( \frac{2mK^2}{n^2\hbar^2} \right)$$

$$E_n = -\frac{mK^2}{2\hbar^2 n^2}$$

$$E_n = -\frac{m}{2\hbar^2}\frac{K^2}{n^2} \hspace{1.5cm} n = p + \ell + 1 = \{1, 2, 3, \cdots \}$$

\subsection{c) Obtenga las primeras funciones del átomo de Hidrógeno}

$$\begin{aligned}
\psi_{100} &= \frac{1}{\sqrt{\pi}a^{3/2}}e^{-\rho} \\
\psi_{200} &= \frac{1}{\sqrt{32\pi}a^{3/2}}(2-\rho)e^{-\rho/2} \\
\psi_{210} &= \frac{1}{\sqrt{32\pi}a^{3/2}}\rho e^{-\rho/2}\cos\theta \\
\psi_{211} &= \frac{1}{\sqrt{64\pi}a^{3/2}}\rho e^{-\rho/2}\sin\theta e^{-i\varphi} \\
\dots 
\end{aligned}$$

\textbf{Tenga en cuenta las siguientes expresiones:}

\textbf{Polinomios Asociados de Laguerre}

Los polinomios de Laguerre generalizados (o asociados), $L_q^p(x)$, se pueden definir a partir de la fórmula de Rodrigues:

$$L_q^p(x) = \frac{x^{-p} e^x}{q!} \frac{\mathrm{d}^q}{\mathrm{d}x^q} \left( e^{-x} x^{q+p} \right)$$

Para el caso de la función de onda radial del átomo de hidrógeno, los índices toman los valores $q = n-\ell-1$, $p = 2\ell+1$ y la variable es $\rho$. Su forma explícita expresada como una serie de potencias es:

$$L_{n-\ell-1}^{2\ell+1}(\rho) = \sum_{k=0}^{n-\ell-1} (-1)^k \frac{(n+\ell)!}{(n-\ell-1-k)! \, (2\ell+1+k)! \, k!} \rho^k$$


\textbf{Armónicos Esféricos}

Los armónicos esféricos, $Y_{\ell m}(\theta, \varphi)$, constituyen la solución a la parte angular de la ecuación y se definen típicamente (incluyendo la convención de fase de Condon-Shortley) como:

$$Y_{\ell m}(\theta, \varphi) = (-1)^m \sqrt{\frac{2\ell+1}{4\pi} \frac{(\ell-m)!}{(\ell+m)!}} P_\ell^m(\cos\theta) e^{im\varphi}$$

Donde esta expresión es válida para $m \geq 0$. Para valores negativos de $m$, se utiliza la relación de conjugación compleja: $Y_{\ell, -m}(\theta, \varphi) = (-1)^m Y_{\ell m}^*(\theta, \varphi)$.

El término $P_\ell^m(\cos\theta)$ corresponde a los **polinomios asociados de Legendre**, los cuales se evalúan en $x = \cos\theta$ y se definen mediante:

$$P_\ell^m(x) = \frac{(-1)^m}{2^\ell \ell!} (1-x^2)^{m/2} \frac{\mathrm{d}^{\ell+m}}{\mathrm{d}x^{\ell+m}} (x^2-1)^\ell$$

\subsubsection{Solución}

Función de onda del estado fundamental $\psi_{100}$, $n=1$, $\ell=0$ y $m=0$.

La ecuación general que tenemos es:

$$\psi_{n\ell m}(r, \theta, \varphi) = \sqrt{ \left( \frac{2}{na_0} \right)^3 \frac{(n-\ell-1)!}{2n(n+\ell)!} } e^{-\rho_{eq}/2} \rho_{eq}^\ell L_{n-\ell-1}^{2\ell+1}(\rho_{eq}) Y_{\ell m}(\theta, \varphi)$$

Cálculos:

\textbf{Parámetro de distancia}

Para $n=1$, la variable de la ecuación se reduce a:

$$\rho_{eq} = \frac{2r}{(1)a_0} = \frac{2r}{a_0}$$

Por lo tanto, el término exponencial será $e^{-\rho_{eq}/2} = e^{-r/a_0}$.

\textbf{Constante de Normalización Radial}

Sustituimos $n=1$ y $\ell=0$ en el término dentro de la raíz:

$$\sqrt{ \left( \frac{2}{(1)a_0} \right)^3 \frac{(1-0-1)!}{2(1)(1+0)!} } = \sqrt{ \frac{8}{a_0^3} \frac{0!}{2(1!)} }$$

Dado que $0! = 1$ y $1! = 1$, simplificamos:

$$\sqrt{ \frac{8}{a_0^3} \left(\frac{1}{2}\right) } = \sqrt{ \frac{4}{a_0^3} } = \frac{2}{a_0^{3/2}}$$

\textbf{Término de potencia y Polinomio de Laguerre}

El término de potencia radial es simplemente $\rho_{eq}^\ell = \rho_{eq}^0 = 1$.

Para el polinomio de Laguerre $L_{n-\ell-1}^{2\ell+1}(\rho_{eq})$, evaluamos los índices:

Índice superior: $2\ell+1 = 2(0)+1 = 1$

Índice inferior: $n-\ell-1 = 1-0-1 = 0$

Necesitamos calcular $L_0^1$. Por definición, el polinomio de Laguerre de grado 0 es siempre 1 (independientemente del superíndice y del argumento):


$$L_0^1(\rho_{eq}) = 1$$

\textbf{Armónico Esférico}

Para $\ell=0$ y $m=0$, calculamos el armónico esférico $Y_{00}$. Sabiendo que el polinomio de Legendre asociado $P_0^0(\cos\theta) = 1$:

$$Y_{00}(\theta, \varphi) = \sqrt{\frac{2(0)+1}{4\pi} \frac{(0-0)!}{(0+0)!}} P_0^0(\cos\theta) e^{i(0)\varphi} = \sqrt{\frac{1}{4\pi}} = \frac{1}{2\sqrt{\pi}}$$

\textbf{Ensamblaje Final}

Ahora multiplicamos todos los componentes que hemos calculado:

$$\psi_{100} = \left( \frac{2}{a_0^{3/2}} \right) e^{-r/a_0} (1) (1) \left( \frac{1}{2\sqrt{\pi}} \right)$$


$$\psi_{100} = \frac{1}{\sqrt{\pi} a_0^{3/2}} e^{-r/a_0}$$



\subsection{d)} Estudiar la densidad radial de probabilidad electrónica. Muestre que la distancia más probable de encontrar el electrón en el estado $(100)$ es justamente el radio de Bohr.

\subsection{e)} La función de onda del estado base es $\psi_{100} = \frac{1}{\sqrt{\pi}a_0^{3/2}}e^{-\rho}$, ¿cuál es la probabilidad de encontrar al electrón a una distancia menor al radio de Bohr?




\bibliographystyle{plainurl}
\bibliography{bibliografia} 
\end{document}
